\documentclass[11pt]{article}

\usepackage[margin=1in]{geometry}
\usepackage{parskip}
\usepackage{amsmath, amssymb, amsthm}
\usepackage{booktabs}
\usepackage{hyperref}

\title{A small library for first-order initial value problems}
\author{Group N}
\date{}

\begin{document}
\maketitle

\begin{abstract}
We present \texttt{ode\_flow}, a small Python library for solving and
sampling first-order initial value problems of the form
\(y'(t) = f(t, y(t))\) with \(y(t_0) = y_0\). The library is a thin
wrapper around \texttt{sympy.dsolve}~\cite{meurer2017}, with
additional helpers to evaluate the solution at a list of times and to
verify that a candidate expression satisfies the ODE. We illustrate
the library on three standard problems: exponential growth, constant
growth, and quadratic growth. The library follows the Diataxis
framework for documentation~\cite{procida2017} and the differential
equations chapter of the open-source textbook \emph{Python for
Mathematics}~\cite{knight2024}.
\end{abstract}

\section{Introduction}

A first-order ordinary differential equation has the form

\[
    y'(t) = f(t, y(t)),
\]

and an initial value problem adds the condition \(y(t_0) = y_0\). For
many right-hand sides \(f\) the solution can be written in closed
form by recognising one of the standard cases (separable, linear,
exact)~\cite{boyce2017}. The library \texttt{sympy}~\cite{meurer2017}
implements these cases through \texttt{sympy.dsolve}, so we treat the
solver as the engine and focus on a small wrapper that makes the
common operations easy.

We present a three-function library,
\texttt{ode\_flow}, that solves a first-order initial value problem,
samples the solution at a list of times, and verifies that a
candidate expression satisfies the ODE.

\section{The library}

The library lives in the single module \texttt{ode\_flow.py} and
follows the differential equations chapter of \emph{Python for
Mathematics}~\cite{knight2024}: we use \texttt{sym.Function} for the
unknown, \texttt{sym.Eq} for the equation, and \texttt{sym.dsolve}
with the \texttt{ics} keyword for the initial condition.

\texttt{solve\_initial\_value\_problem(right\_hand\_side,
time\_variable, function, initial\_time, initial\_value)} builds the
equation, calls \texttt{sym.dsolve}, and returns the right-hand side
of the resulting equation \texttt{y(t) = expression}.

\texttt{trajectory(solution\_expression, time\_variable,
sample\_times)} evaluates the solution at a list of times. Each
evaluation is cast to \texttt{float} so the output is easy to plot.

\texttt{verify\_solution(solution\_expression, right\_hand\_side,
time\_variable, function)} substitutes the candidate into the
right-hand side and compares with the derivative of the candidate;
the two simplify to the same expression for a genuine solution.

\section{Worked examples}

The problem \(y' = y\) with
\(y(0) = 1\) has the well-known solution \(y(t) = e^t\). The library
recovers this exactly; sampling at \(t = 0, 1, 2\) gives
\(1, e, e^2\) cast to floats. The verifier accepts \(e^t\) and
rejects \(t\).

The problem \(y' = 1\) with \(y(0) = 0\)
has the solution \(y(t) = t\). The library returns this in one line.

The problem \(y' = 2t\) with \(y(0) = 0\)
has the solution \(y(t) = t^2\). Sampling at the integer times
\(0, 1, 2, 3\) gives the squares \(0, 1, 4, 9\).

The problem \(y' = y\) with
\(y(1) = e\) also has the solution \(y(t) = e^t\); the verifier
accepts it. This is useful because it confirms that the initial
condition does not need to be at \(t_0 = 0\).

\section{Discussion}

\texttt{sym.dsolve} can already solve
everything in this paper. The library exists because we wanted a
shorter interface for the common case (first-order, scalar, initial
value problem), and because building the equation and the
\texttt{ics} dictionary by hand each time is fiddly enough that we
preferred to encapsulate it.

The library handles only scalar first-order
problems. Systems of ODEs, higher-order equations, and boundary value
problems are all natural extensions. A natural next step is to allow
the right-hand side to depend on extra parameters and to expose the
sweep over those parameters as a separate function.

\section{Conclusion}

We have presented \texttt{ode\_flow}, a three-function library for
first-order initial value problems. The library recovers the
closed-form solutions of three standard problems and confirms them
through both sampling and symbolic verification.

\bibliographystyle{plain}
\bibliography{references}

\end{document}
